\section{Technical overview}\label{sec:overview}

\subsection{From a univariate oracle to a multivariate degree}\label{sec:ov-idea}

The prover commits to $f\in\F[x_1,\dots,x_m]_{<k}$ through the evaluations of its
Kronecker image $U_f(X)=f(X,X^k,\dots,X^{k^{m-1}})$ on a smooth domain $L_0$. FRI run on
this word certifies proximity to $\RS[\F,L_0,k^m]$, and by \cref{lem:kron} this says only
that every individual degree of $f$ is below $k$. The total degree is invisible to it: the
monomial $x_1^{k-1}\cdots x_m^{k-1}$ and the linear polynomial $x_1$ both have Kronecker
images of degree below $k^m$.

FRI nevertheless comes close to seeing $f$. A $k$-ary fold with a scalar challenge $r_i$
replaces the variable $x_i$ by $r_i$ (\cref{lem:fold-sub}): the folded word is the Kronecker
image of $f(r_1,\dots,r_i,x_{i+1},\dots,x_m)$, and after $m$ rounds the last value is
$f(r_1,\dots,r_m)$. For multilinear $f$ and $k=2$ this is the observation behind
BaseFold~\cite{ZCF24}. A single field element, however, carries no information about the
degree of $f$.

Our protocol keeps the substitution and makes the challenge symbolic. In round $i$ the
verifier sends $r_i\in\F$, and the fold uses the challenge $r_iT^{w_i}$, where $T$ is a formal
variable and $w\in\N^m$ is a fixed weight vector:
\[
x_i\ \longmapsto\ r_iT^{w_i}\quad(i=1,\dots,m),\qquad
g(T)=f\big(r_1T^{w_1},\dots,r_mT^{w_m}\big)=\sum_\alpha f_\alpha\,r^\alpha\,T^{\langle w,\alpha\rangle},
\]
where $r:=(r_1,\dots,r_m)$, $r^\alpha:=r_1^{\alpha_1}\cdots r_m^{\alpha_m}$, and
$f_\alpha$, $\langle w,\alpha\rangle$ are as in \cref{sec:notation}.
The last message is the polynomial $g$, and its degree is the weighted degree $\deg_wf$
unless $f^{\mathrm{top}}(r)=0$, where
$f^{\mathrm{top}}:=\sum_{\langle w,\alpha\rangle=\deg_wf}f_\alpha x^\alpha$ collects the terms of
largest weighted degree; this happens with probability at most
$m(k-1)/|\F|$ (\cref{prop:complete}). The verifier checks $\deg g\le D$. With
$w=(1,\dots,1)$ this tests the total degree; with $w=e_l$ it tests the degree in $x_l$.

The price is in the messages. The intermediate oracles take values in $\F[T]$, and an honest
value at level $i$ has degree at most $d_i=\min\big(D,(k-1)\sum_{l\le i}w_l\big)$
(\cref{lem:stab}). For $w=(1,\dots,1)$ a value is a vector of at most $D+1$ field elements,
and the verifier performs the same fibre computations as in FRI, with coefficients in $\F[T]$
instead of $\F$. For $k=4$, $w=(1,\dots,1)$ and $D=6$, the caps are $d_1=3$ and
$d_i=6$ for $i\ge2$. \Cref{fig:chain} shows the polynomial behind each oracle.

\begin{figure}[t]
\centering
\begin{tikzpicture}[node distance=0.55cm,>=stealth,every node/.style={font=\small}]
\node (a) {$L_0:\quad f(x_1,x_2,\dots,x_m)$};
\node (b) [below=of a] {$L_1:\quad f(r_1T^{w_1},x_2,\dots,x_m)$};
\node (c) [below=of b] {$\vdots$};
\node (d) [below=of c] {$L_{m-1}:\quad f(r_1T^{w_1},\dots,r_{m-1}T^{w_{m-1}},x_m)$};
\node (e) [below=of d] {$g(T)=f(r_1T^{w_1},\dots,r_mT^{w_m}),\quad \text{check }\deg g\le D$};
\draw[->] (a) -- node[right]{\scriptsize$\Fold_{r_1T^{w_1}}$} (b);
\draw[->] (b) -- (c);
\draw[->] (c) -- (d);
\draw[->] (d) -- node[right]{\scriptsize$\Fold_{r_mT^{w_m}}$} (e);
\end{tikzpicture}
\caption{The honest oracles. The word on $L_i$ is the Kronecker image of the polynomial shown,
evaluated on $L_i$; its values are polynomials in $T$ of degree at most $d_i$.}
\label{fig:chain}
\end{figure}

\subsection{What has to be proved}\label{sec:ov-target}

Completeness follows from the fact that a fold is a substitution; the one point that needs
care is the $T$-degree cap, which \cref{lem:stab}(c) provides. Soundness asks for more than
FRI's. The verifier must reject, with high probability, any committed word that is far from
the subcode $C_{w,D}\subseteq\RS[\F,L_0,k^m]$ of Kronecker images of polynomials with
$\deg_wf\le D$. Along the way, the honest words belong to the graded codes
$\Gc_i^{w,D}$ (Definition~\ref{def:graded}): their values lie in $\F[T]$, and only the words
coming from a polynomial of graded degree at most $D$ are allowed.

The soundness analysis of FRI by Ben-Sasson, Carmon, Ishai, Kopparty and
Saraf~\cite[\S8]{BCIKS20} reaches the Johnson radius, and we follow it. It applies to words
with values in $\F$ and to Reed--Solomon codes, and both assumptions fail here: our words take
values in $\F[T]$, and our codes are proper subcodes. We remove each obstacle with one
observation, and neither requires a new decoding argument.

\subsection{Two observations}\label{sec:ov-ideas}

\paragraph{Polynomial values as extension-field values.} Fix $t$ larger than every
$T$-degree that occurs, and identify $\F[T]_{<t}$ with a degree-$t$ extension $\E$ of $\F$
through an $\F$-basis. A word with values in $\F[T]_{<t}$ whose coefficients are
Reed--Solomon codewords becomes a Reed--Solomon codeword over $\E$ on the same domain, and the
identification preserves agreement and $\F$-linear combinations (\cref{lem:ext}). One round
of folding is then a curve
\[
\Fold_{rT^{w_i}}(u)=\sum_{j<k}r^j\,\big(T^{w_ij}u_j\big)
\]
with coefficients over $\E$ and parameter $r\in\F$. The theorems of~\cite{BCIKS20} draw the
parameter from the field of the code, and their Remark~1.2 allows a larger parameter field,
not a smaller one. Their proofs, however, use the set of good parameters only through its size
and through membership in it (\cref{ext:proof-facts}, with page references), so they apply
as they stand, with the error measured in $|\F|$ (\cref{prop:H1}). Of the two facts we use about these proofs, one is a consequence of a stated proposition
of~\cite{BCIKS20}, and the other is a reading of the proofs: it is the step of the paper that
does not rest on a stated theorem, and \cref{sec:params} records it with the places where it
can be checked.

\paragraph{Correlated agreement passes to linear subcodes.} Let $W$ be an $\F$-linear
subcode. For each good parameter $r$ we choose the nearby codeword $P_r$ inside $W$. The
proof of~\cite{BCIKS20} produces codewords $v_0,\dots,v_l$ with $\sum_jr^jv_j=P_r$ for more
than $l$ values of $r$ (\cref{ext:proof-facts}, (F1)). Choosing $l+1$ of them,
\[
\begin{pmatrix}P_{r_0}\\ \vdots\\ P_{r_l}\end{pmatrix}
=\big(r_a^{\,j}\big)_{0\le a,j\le l}
\begin{pmatrix}v_0\\ \vdots\\ v_l\end{pmatrix},
\]
and the Vandermonde matrix is invertible over $\F$, so each $v_j$ is an $\F$-linear
combination of the $P_r$ and lies in $W$. Weighted correlated agreement over curves is
therefore inherited by every $\F$-linear subcode, with the same parameters
(\cref{cor:share}).

\paragraph{Assembly.} In the round lemma (\cref{lem:round}) the two observations give
codewords $v_j$ of the graded code that agree with the curve coefficients $T^{w_ij}u_j$ on a
set of large weight. Agreement on a large set forces each $v_j$ to be divisible by
$T^{w_ij}$, and unfolding (\cref{lem:unfold}) turns the quotients into a word of the
previous graded code that agrees with the previous oracle on all fibres above that set. The
proof of \cref{thm:main} is then the weight argument of~\cite[Lemma~8.2]{BCIKS20}, run along
the chain $\Gc_0^{w,D},\dots,\Gc_m^{w,D}$. The last code consists of the constant words with
value in $\F[T]_{\le D}$, which is exactly what the verifier checks on $g$.

\subsection{Scope}\label{sec:ov-scope}

\Cref{cor:share} uses nothing about degrees: it applies to any $\F$-linear subcode, for
instance to the constrained Reed--Solomon codes of WHIR~\cite{ACFY25}. Our soundness theorem
is accordingly a reduction. It needs weighted correlated agreement over curves for
Reed--Solomon codes, with the two properties of the proof recorded in \cref{sec:params};
any such theorem beyond the Johnson radius would carry over to the degree test without
change.
